% Vietnamese translation for OI Public Library
% Source-File: IOI中国国家候选队论文/2005/吴景岳/吴景岳.ppt
\documentclass[11pt]{article}
\usepackage{oipl-vn}

\title{Bản trình chiếu: Ghi chú về tối ưu trên lưới}
\author{Wu Jingyue}
\date{2005}

\begin{document}
\maketitle

\origtitle{Slide companion for Wu Jingyue's grid presentation}
\sourcepath{IOI China candidate team papers / 2005 / Wu Jingyue / Wu Jingyue.ppt}

\section{Phạm vi bản dịch}

Tệp gốc chỉ có PowerPoint. Phần văn bản đọc được gồm mô hình lưới
\(M\times N\), gốc \((0,0)\), ràng buộc \(M,N\le 4000\), các bước
\((r+P,c+Q)\), \((r+Q,c+P)\), \(P,Q\le 10\), \(B\le 10\), và \(D\le 5000\).

\section{Mô hình nhận diện được}

Các slide mô tả một bài toán di chuyển hoặc phủ trên lưới. Từ một ô
\((r,c)\), trạng thái có thể chuyển sang hai vị trí đối xứng theo \(P,Q\).
Khi \(P,Q\) nhỏ nhưng \(M,N\) lớn, cách tiếp cận tự nhiên là dùng BFS, quy
hoạch động theo lớp, hoặc nén trạng thái theo phần dư.

\section{Thông điệp}

Với lưới lớn, không thể lưu mọi biến thể nếu mỗi ô còn kèm nhiều tham số. Cần
khai thác giới hạn nhỏ của \(P,Q,B\) hoặc \(D\) để giảm trạng thái. Bản dịch
này giữ lại các ràng buộc và mô hình chuyển trạng thái chắc chắn đọc được từ
slide.

\end{document}
